% Обзор литературы 0005.
% Сборка: xelatex review.tex && bibtex review && xelatex review.tex && xelatex review.tex
% Либо: tectonic review.tex
% PDF кладётся в out/, в репозиторий не коммитится.
\documentclass[11pt,a4paper]{article}

\usepackage{iftex}
\ifPDFTeX
  \usepackage[T2A]{fontenc}
  \usepackage[utf8]{inputenc}
  \usepackage[russian]{babel}
\else
  \usepackage{fontspec}
  \usepackage[russian]{babel}
  \IfFontExistsTF{Times New Roman}
    {\setmainfont{Times New Roman}}
    {\IfFontExistsTF{DejaVu Serif}
      {\setmainfont{DejaVu Serif}}
      {\IfFontExistsTF{Liberation Serif}
        {\setmainfont{Liberation Serif}}
        {}}}
\fi

\usepackage[margin=2.4cm]{geometry}
\usepackage{amsmath}
\usepackage{booktabs}
\usepackage{longtable}
\usepackage{array}
\usepackage{enumitem}
\usepackage{hyperref}
\usepackage{microtype}

\setlength{\emergencystretch}{2em}

\hypersetup{
  colorlinks=true,
  linkcolor=black,
  citecolor=black,
  urlcolor=blue
}

\setlist[itemize]{leftmargin=1.4em,itemsep=0.2em,topsep=0.2em}
\setlist[enumerate]{leftmargin=1.6em,itemsep=0.2em,topsep=0.2em}

\newcommand{\code}[1]{\texttt{#1}}

\title{Снижение токенов на решённую задачу:\\
обзор литературы}
\author{Обзор 0005}
\date{8 ноября 2026}

\begin{document}
\maketitle

\begin{abstract}
\noindent
Обзор отвечает на вопрос, как снизить расход токенов на одну решённую задачу
у кодового агента при фиксированном бюджете, не теряя долю решённых задач.
Включено 64 работы: измерение стоимости на задачу, маршрутизация моделей и
режимов рассуждения, ограничение длины рассуждения, сжатие контекста,
управление траекторией, сокращение числа шагов, переиспользование планов,
ранний отказ и распределение бюджета. Вывод: в метрике
$C=O+\lceil I/10\rceil$ выходные токены весят в десять раз больше входных,
поэтому наибольший рычаг дают выбор режима рассуждения и длины рассуждения,
а не сжатие промпта. Второй по важности рычаг — число шагов: смена интерфейса
действий и отделение исследования репозитория от решения дают 40--60\,\%
экономии. Третий — ранний отказ от бесперспективных задач: неудачная задача
тоже стоит токенов. При этом сокращение токенов не равно сокращению
стоимости: сжатие меняет траекторию, добавляет шаги и может ухудшить успех.
Ни одна работа не измеряет нашу метрику на четырёх моделях Qwen с общим
бюджетом, поэтому вывод обзора — не спецификация, а программа
эксперимента.
\end{abstract}

\section{Вопрос и зачем}

\subsection{Вопрос}

Как снизить расход токенов на одну решённую задачу у кодового агента с
инструментами при фиксированном бюджете, не теряя долю решённых задач?

Регламент считает стоимость вызова в выходных токенах:
\begin{equation}
  C \;=\; O + \left\lceil \frac{I}{10} \right\rceil ,
  \label{eq:cost}
\end{equation}
где $O$ — выходные токены, включая рассуждения, а $I$ — входные. Кеш
провайдера не учитывается: все входные токены считаются полностью. Баллы
начисляются только за задачи, прошедшие скрытые тесты. Поэтому целевая
величина — не средняя стоимость задачи, а
\begin{equation}
  \text{tokens per solved task}
  \;=\; \frac{\sum_{\text{все задачи}} C}{\#\{\text{решённые задачи}\}} .
  \label{eq:metric}
\end{equation}

Формула \eqref{eq:metric} раскладывается на четыре независимых рычага:

\begin{enumerate}
  \item \textbf{Выход на шаг} — длина рассуждения, формат ответа, число
        пробных попыток внутри одного вызова.
  \item \textbf{Вход на шаг} — размер контекста, истории, вывода инструментов
        и системной инструкции. Вход в десять раз дешевле выхода.
  \item \textbf{Число шагов} — сколько вызовов модели и обращений к
        инструментам нужно, чтобы прийти к ответу.
  \item \textbf{Число неудачных задач} — задачи, которые не решены, всё
        равно увеличивают числитель и не увеличивают знаменатель.
\end{enumerate}

К этим рычагам добавляется распределение бюджета: какие задачи решать
первыми и когда прекращать прогон.

PICO:

\begin{itemize}
  \item \textbf{Population} — кодовые агентские задачи: исправление issue в
        репозитории, генерация кода, терминальные задачи, а также задачи на
        рассуждение, если эффект переносится на агентский цикл. Один агент
        с инструментами, цикл ReAct или его вариант, общий бюджет токенов.
  \item \textbf{Intervention} — приём, меняющий расход: выбор модели и
        режима рассуждения на шаг, ограничение длины и ранний выход, сжатие
        контекста и вывода инструментов, управление траекторией и памятью,
        сокращение числа шагов и смена интерфейса действий, ранний отказ.
  \item \textbf{Comparison} — тот же агент без приёма: полное рассуждение,
        полный контекст, самая сильная модель на каждом шаге, максимальное
        число попыток. Либо другой приём при том же бюджете.
  \item \textbf{Outcome} — стоимость \eqref{eq:cost} на решённую задачу,
        доля решённых задач, разброс стоимости, стоимость одной решённой
        задачи.
\end{itemize}

\subsection{Зачем}

Харнесс получает 100 задач, общий бюджет и один час. Доступны четыре
модели, включая варианты \code{-noreason}. Харнесс сам выбирает модель и
может переключать её на любом вызове. Это значит, что решения о расходе
принимаются не один раз, а на каждом шаге.

Ответ на вопрос определяет:

\begin{itemize}
  \item какие механизмы ставить первыми: маршрутизацию режима рассуждения,
        ограничение длины, управление контекстом, смену интерфейса действий
        или ранний отказ;
  \item на что тратить ограниченное время разработки до прогона;
  \item какой эксперимент поставить и какие метрики в нём мерить;
  \item нужна ли спецификация изменения поведения, если приём доказал
        прирост.
\end{itemize}

Обзор не проектирует харнесс и не выбирает конкретные пороги. Он собирает,
что уже измерено, при каких условиях эффект воспроизводится и где остаётся
пробел.

\section{Метод}

\subsection{Базы и запросы}

Поиск выполнен 8 ноября 2026 года в двух базах: arXiv API (поле
\code{abs:}) и Crossref (поле \code{query.bibliographic}). Просматривались
верхние 20 записей на запрос, сортировка по релевантности. Язык —
английский, годы — 2015--2026. Выполнено 42 запроса к arXiv и 6 запросов к
Crossref; два запроса Crossref вернули HTTP 429 и не дали записей.

Semantic Scholar и OpenAlex не использовались: в прошлом обзоре они вернули
HTTP 403 и дневной лимит. Crossref плохо ранжирует CS-препринты, поэтому
почти все включённые работы найдены в arXiv. Точные запросы, даты, числа
найденных и просмотренных записей приведены в \code{protocol.md}.

У Crossref запрос без кавычек возвращает всю индексную выдачу, поэтому
\code{total-results} в миллионы записей не отражает число релевантных
работ. Для отбора считались только верхние 20 записей на запрос.

\subsection{Критерии отбора}

Включались работы, которые:

\begin{itemize}
  \item предлагают или оценивают приём, меняющий расход токенов, длины
        рассуждения, числа шагов или числа попыток;
  \item дают количественный результат рядом с метрикой качества или успеха;
  \item относятся к коду, software engineering, терминальным или агентским
        задачам; задачи на рассуждение допускались, если приём переносится
        на агентский цикл;
  \item доступны по ссылке.
\end{itemize}

Исключались работы, где ускорение не меняет число токенов (квантизация,
speculative decoding, KV-кеш, батчинг), где кеш провайдера — единственный
механизм экономии, где нет замера на инференсе, где нет метрики качества
рядом с расходом и где домен не переносится.

\subsection{Скрининг и извлечение}

Один просматривающий (coding-агент). Сначала заголовки и аннотации: работа
отбрасывалась, если она не про расход токенов или не про агентский цикл.
Затем аннотация или полный текст для отобранных: проверялось наличие
количественного результата и метрики качества. Для каждой включённой работы
выписаны ссылка, год, подход, данные и домен, метрика, результат и
ограничение. Числа извлечены из аннотаций; полные тексты проверялись
выборочно.

\subsection{Оценка качества}

Признаки надёжности: рецензируемость, открытость кода и данных, число задач
и моделей, контролируемое сравнение с базовой линией, наличие метрики
успеха рядом с метрикой расхода, воспроизводимость и совпадение условий с
нашим доменом: кодовый агент, API-модели, фиксированный бюджет.

\subsection{Синтез}

Нарративный разбор по группам приёмов плюс сводная таблица. Мета-анализ не
проводится: работы меряют разные величины (длина CoT, входные токены,
доллары, FLOPs, число шагов), а домены и модели несопоставимы.

\subsection{Подсчёт отбора (PRISMA)}

\begin{table}[h]
\centering
\begin{tabular}{lr}
\toprule
Этап & Записей \\
\midrule
Идентифицировано в базах & 4\,511 \\
Просмотрено верхних записей (20 на запрос) & 822 \\
Удалено дубликатов & 23 \\
Просмотрено заголовков и аннотаций & 799 \\
Исключено на скрининге & 686 \\
Запрошено аннотаций и полных текстов & 113 \\
Исключено после полного текста & 49 \\
Включено в обзор & 64 \\
\bottomrule
\end{tabular}
\caption{Схема отбора PRISMA.}
\end{table}

Для arXiv «идентифицировано» — сумма \code{totalResults} по 42 запросам
(4\,431 запись); для Crossref — верхние 20 записей на каждый из четырёх
успешных запросов (80 записей). Всего просмотрено 742 записи arXiv и 80
записей Crossref. Дубликаты: 22 внутренних повтора arXiv и одно пересечение
между базами. После скрининга осталось 113 работ, у которых прочитана
аннотация или полный текст; 49 исключены с причиной, 64 включены.

\section{Что известно}

\subsection{Мерить надо стоимость на решённую задачу, а не токены}

Первая группа работ задаёт метрику и показывает, почему сырые токены плохо
предсказывают результат.

Cost-of-Pass \cite{erol2025costpass} формализует стоимость правильного
ответа как ожидаемые денежные затраты и строит границу минимальной
стоимости по доступным моделям. Авторы показывают, что лёгкие модели
выгодны для базовых количественных задач, крупные — для задач на знания,
reasoning-модели — для сложных; majority voting и self-refinement редко
оправдывают свою цену, а бюджетный TALE-EP выглядит перспективно. Та же
идея в агентском исполнении: Efficient Agents \cite{wang2025efficientagents}
сохраняет 96{,}7\,\% качества OWL на GAIA при цене \$0{,}228 вместо
\$0{,}398, то есть улучшает cost-of-pass на 28{,}4\,\%.

Измерения на кодовых агентах дают масштаб проблемы. Bai et al.
\cite{bai2026aiagents} анализируют траектории восьми моделей на SWE-bench
Verified: агентские задачи тратят примерно в 1000 раз больше токенов, чем
code chat; входные токены доминируют в сыром счёте; один и тот же прогон
может отличаться до 30 раз; точность растёт до среднего расхода и выходит
на плато, а не растёт монотонно. Tokenomics \cite{salim2026tokenomicsquantify}
разлагает расход ChatDev по фазам: code review — 59{,}4\,\% токенов,
входные токены — 53{,}9\,\%. Основной расход не в первичной генерации, а в
доработке и проверке.

Отсюда два следствия. Во-первых, средняя стоимость задачи без учёта успеха
обманчива: дорогая неудача и дешёвая победа входят в одну метрику
по-разному. Во-вторых, сокращение токенов не равно сокращению стоимости.
Weinberger и Hozez \cite{weinberger2026tokenreduction} сравнивают три
сжатия вывода инструментов с немодифицированным Claude Code: самое
агрессивное сжатие убрало 38{,}4\,\% токенов вывода, но увеличило
оплаченную стоимость на 6{,}8\,\%; корреляция экономии токенов и стоимости
составила $r=0{,}15$. Сжатие меняло траекторию — агент чаще обращался к
источникам, диагностировал и тестировал, — а на SWE-bench Go агрессивное
сжатие снизило успешность применения патча.

Приёмы оркестрации дают больший эффект, чем микрооптимизации промпта.
Смена harness при фиксированных моделях \cite{ali2026harnesseffect} даёт
$-38\,\%$ токенов на задачу, $-41\,\%$ стоимости и качество на уровне;
эффект сохраняется для всех шести моделей. Local-Splitter
\cite{agyemang2026localsplitter} измеряет семь тактик на четырёх классах
кодовых нагрузок: локальная маршрутизация плюс сжатие промпта дают
45--79\,\% экономии облачных токенов на edit- и explanation-нагрузках, а
на RAG-нагрузке полный набор с draft-review — 51\,\%. TokenEconomizer
\cite{li2026tokeneconomizerexp} снижает стоимость на 18{,}4--30{,}9\,\% при
сопоставимом успехе за счёт микро-навыков, удаления нерелевантных
инструментов и вмешательства в непродуктивные циклы. OckBench
\cite{du2025ockbenchmeasuring} фиксирует разброс: модели с близкой
точностью могут различаться по длине решения до пяти раз.

Практический вывод для нас: любую оптимизацию надо мерить на выходе
\eqref{eq:metric}, с успехом и распределением, а не по числу сэкономленных
токенов. Иначе легко получить $-38\,\%$ токенов и $+6{,}8\,\%$ стоимости.

\subsection{Маршрутизация модели и режима рассуждения}

Если выходные токены весят в десять раз больше входных, самый прямой рычаг —
не генерировать их там, где они не нужны. Литература о каскадах и
маршрутизации даёт самые большие числа в обзоре.

FrugalGPT \cite{chen2023frugalgptuse} формулирует три стратегии экономии:
адаптацию промпта, приближение модели и каскад. Каскад, который учится
выбирать комбинацию моделей под запрос, достигает качества GPT-4 при
снижении стоимости на 98\,\% или улучшает качество на 4\,\% при той же
цене. Каскад по согласию ответов \cite{yue2023largelanguage} даёт качество
GPT-4 за 40\,\% его стоимости. Для кода каскад с self-testing
\cite{chen2024modelcascading} снижает стоимость в среднем на 26\,\%, а в
лучшем случае на 70\,\%, сохраняя или улучшая точность. Пошаговая
маршрутизация STEER \cite{lee2025confidenceguided} выбирает модель на
каждом шаге по уверенности и на AIME даёт до $+20\,\%$ точности при
$-48\,\%$ FLOPs против большой модели.

Для агентов важнее не выбор модели на задачу, а выбор на шаг.
Budget-Aware Agentic Routing \cite{zhang2026budgetaware} выбирает дешёвую
или дорогую модель на каждом шаге и улучшает границу cost--success при
жёстком бюджете. Ares \cite{yang2026aresadaptive} предсказывает минимальный
уровень рассуждения для шага и снижает reasoning-токены до 52{,}7\,\% при
минимальной потере успеха. DART \cite{lee2026dartdraft} обходится без
обучения: генерирует два дешёвых черновика без рассуждения, принимает
прямой ответ при их согласии и оценивает бюджет рассуждения по энтропии при
расхождении. На математике и коде это снижает thinking-токены на
32--73\,\% и иногда повышает точность. HRBench \cite{ning2026hrbenchbenchmarkin}
сравнивает семейства переключения режима и находит, что prompt-based выбор
даёт хороший trade-off точности и токенов, внешний routing — более
стабильную экономию, а speculative-исполнение обычно дороже.

Отдельно стоит ранний отказ. Каскад с отказом на ранних моделях
\cite{zellinger2025costsaving} снижает стоимость на 13\,\% и ошибки на
5\,\% ценой роста доли отказов на 4{,}1\,\%. Это прямо связано с
невыполнимыми задачами: если задача не будет решена, лучше не доводить её
до дорогой модели. HERA \cite{luo2026heraharness} доводит abstention
accuracy с 61{,}7\,\% до 83{,}3\,\% и переносит навык на 19 моделей,
снижая стоимость малых моделей примерно на 85\,\%.

Ограничение группы: каскады и роутеры обучены на других моделях и
распределениях. Для четырёх Qwen-моделей порог переключения придётся
калибровать заново. Но сама доступность вариантов \code{-noreason} делает
маршрутизацию самыми дешёвыми в реализации и потенциально самыми
эффективными воротами экономии.

\subsection{Ограничение длины рассуждения}

Вторая группа работ сокращает выходные токены внутри одного вызова.

TALE \cite{han2024tokenbudget} добавляет в промпт бюджет токенов и
подбирает его по сложности задачи; это заметно сокращает CoT при небольшом
падении точности. Token Complexity \cite{lee2025wellllms} объясняет, почему
единый бюджет плох: у каждой задачи есть порог токенов, ниже которого
точность резко падает, а промптовые стратегии сжатия работают далеко от
теоретического предела. Отсюда вывод — сжатие должно быть адаптивным по
трудности.

Адаптивные методы дают 20--50\,\% длины без заметной потери. Dynamic Early
Exit \cite{yang2025dynamicearly} останавливает генерацию по уверенности в
пробном ответе: на 11 моделях и 10 бенчмарках длина падает на
19{,}1--80{,}1\,\%, а точность иногда растёт на 0{,}3--5{,}0\,\%. ESTAR
\cite{wang2026estarearly} обучает модель сигналу остановки и сокращает
длину в 3{,}7 раза (с 4\,799 до 1\,290 токенов) при точности 74{,}9\,\%
против 74{,}2\,\%. DTSR \cite{xiang2026thinkingenough} снижает длину на
28{,}9--34{,}9\,\%. Budget Guidance \cite{li2025steeringllm} без
дообучения удерживает 63\,\% токенов полного рассуждения при конкурентной
точности. L1 \cite{aggarwal2025l1controlling} и s1 \cite{muennighoff2025s1simple}
показывают, что длиной можно управлять грубо и точно, а CARE
\cite{he2026thinkor} добавляет 4\,\% Pass@1 при $-37\,\%$ длины. Confidence
Training \cite{hosseini2026learningstop} снижает токены на 25\,\% при той же
точности, не вводя явного штрафа за длину.

Но простое «думай короче» ломается на трудных задачах. В кодовых
reasoning-моделях \cite{xue2025empiricalstudy} умеренное сокращение шагов
сохраняет успех на стандартных задачах и редко — на трудных. Для агентов
overthinking опаснее: на SWE-bench Verified \cite{cuadron2025dangeroverthinking}
выбор решения с меньшим overthinking даёт почти $+30\,\%$ успеха и
$-43\,\%$ стоимости. Остановка на первом правильном префиксе
\cite{caldarella2026thinkingpast} повышает точность до 21\,\%, но обычный
early stopping убирает «болтливое» overthinking и не убирает вредное.
Маргинальная полезность токенов падает \cite{zhou2026thinkinghurts}, а
оптимальная длина зависит от трудности.

Для нас это означает: cap на рассуждение нужен, но не единый. Разумная
схема — начинать с короткого бюджета, повышать его при признаках
неуверенности или провала тестов и останавливать генерацию, когда ответ уже
стабилен. Проверять это надо на кодовых задачах: большинство работ группы
сделаны на математике.

\subsection{Сжатие контекста и вывода инструментов}

Третья группа сокращает входные токены. В нашей метрике вход в десять раз
дешевле выхода, поэтому прямой эффект здесь ограничен: если выход не
меняется, сжатие входа даёт не более примерно 9\,\% взвешенной стоимости.
Но сжатие может уменьшить число шагов и предотвратить переполнение
контекста, и тогда его эффект больше.

LLMLingua \cite{jiang2023llmlinguacompressi} сжимает промпт до 20 раз с
небольшой потерей качества. Для агентов важнее сжатие истории и вывода
инструментов. FOCUS \cite{dixit2026focustraining} без обучения сжимает
историю, сохраняя причинно значимые для будущих решений фрагменты:
пиковый контекст $-48\,\%$, зависимости $-73\,\%$, успех $+8{,}9$ п.п.
Латентное сжатие вывода инструментов \cite{liu2026toolsneed} снижает
потребление токенов на 70--89\,\% и время на 24--56\,\%, иногда улучшая
качество. DTOC \cite{chaturvedi2026dtocdynamic} делает сжатие обратимым:
для отзывчивых моделей вход $-10{,}3$--$12{,}7\,\%$, шаги $-2{,}4$--$32{,}3\,\%$,
solve rate в 1{,}5--2{,}5 раза выше, cost per solved task в 3--3{,}5 раза
ниже. Для других моделей эффект смешанный — это важное предупреждение.

Сжатие контекста в кодовых агентах чаще выигрывает не само по себе, а как
способ не упереться в окно. В покомпонентном исследовании harness
\cite{fan2026empiricalstudy} управление контекстом становится ценным при
узком окне, и большая часть пользы — предотвращение отказов из-за
переполнения; rule-based elision до суммаризации даёт лучший режим, а
делать удалённое восстановимым часто не окупается. AGENTS.md
\cite{lulla2026impactagents} связан со снижением выходных токенов на
16{,}58\,\% и runtime на 28{,}64\,\% при сопоставимом завершении: дешёвая
входная инструкция экономит дорогие выходные шаги. ContextPipe
\cite{xu2026contextpipedatabas} уменьшает объём токенов на 31\,\% и число
вызовов на 23\,\% за счёт детерминированной сборки контекста.

Обратная сторона: экономия входа может быть съедена лишними шагами.
Weinberger и Hozez \cite{weinberger2026tokenreduction} показывают это
напрямую. Chen et al. \cite{chen2026measurebefore} добавляют, что равный
бюджет токенов не означает равный доставленный контекст: разные типы
объектов рабочей памяти сжимаются по-разному, а выигрыш калибровки не
переносится на отложенные задачи. Поэтому сжатие должно быть обратимым,
селективным и проверенным на траектории, а не на числе токенов.

\subsection{Управление траекторией и памятью}

Четвёртая группа работает с историей агента: что хранить, что сжимать, что
выгружать во внешнюю память.

ACM \cite{li2026acmagentic} даёт агенту инструменты управления контекстом и
внешнюю память: агент сам решает, когда сжимать, а что сохранить. CAT
\cite{liu2025contexttool} превращает контекст в вызываемый инструмент и на
SWE-bench Verified достигает 57{,}6\,\% solved при ограниченном бюджете.
Context Folding \cite{sun2025scalinglong} сворачивает подтраектории: активный
контекст в 10 раз меньше при уровне ReAct или лучше. ICLR
\cite{wang2026agentsforget} показывает, что исторические рассуждения можно
забывать, если состояние уже вынесено в код, файлы или вывод инструментов:
вход $-25{,}5\,\%$, выход $-14{,}4\,\%$, reward растёт. SWE-MeM
\cite{gao2026swemem} обучает агента проактивной памяти и достигает
43{,}4\,\% и 60{,}2\,\% resolve для моделей 4B и 30B. TACO
\cite{ren2026selfevolving} самоэволюционирующее сжимает наблюдения
терминала, TraceLite \cite{xiao2026traceliteefficient} заменяет старые
наблюдения изображениями и удаляет дубли: $-33{,}05\,\%$ входа и
$-21{,}30\,\%$ стоимости API при 78{,}25\,\% resolve. FastContext
\cite{zhang2026fastcontexttrainin} отделяет исследование репозитория от
решения и снижает токены на 60\,\% при росте resolve до 5{,}5 п.п.

Общая схема: свежие наблюдения хранить точно, старые — сжимать или
выгружать, восстановление делать по требованию, а решение о сжатии
привязывать к прогрессу задачи и остатку бюджета. Но реализация важнее
названия: в \cite{fan2026empiricalstudy} лишняя машинерия восстановления
почти не использовалась моделями и не давала точности.

\subsection{Меньше шагов и лучше интерфейс действий}

Пятая группа сокращает не длину одного ответа, а число ответов.

Agentless \cite{xia2024agentlessdemystify} заменяет автономного агента
тремя фазами — локализация, ремонт, валидация — и на SWE-bench Lite
получает 32{,}00\,\% (96 исправлений) при \$0{,}70, лучше открытых агентов
того времени. Это сильный аргумент, что сложный цикл не всегда нужен.
Сравнение архитектур инструментов \cite{xu2026devilinterface} на 11\,700
траекториях показывает: CodeAct-подобный интерфейс даёт тот же результат
при $-41{,}6\,\%$ шагов и $-56{,}3\,\%$ токенов. Дешёвый критик
высокоуровневого плана \cite{gandhi2026steerdon} повышает resolved rate на
16{,}0\,\% и 14{,}4\,\% для двух агентов и снижает стоимость примера с
\$0{,}07 до \$0{,}03, потому что задача решается за меньшее число шагов.
ContextPipe \cite{xu2026contextpipedatabas} и FastContext
\cite{zhang2026fastcontexttrainin} действуют в ту же сторону.

Отдельный вопрос — не является ли агент вообще лишним. На 216 задачах
Codeforces \cite{dong2026independentsamplin} независимое $k$-shot
сэмплирование выигрывает у агентского рассуждения и по точности на
стоимость, и по точности на вызов. Для самодостаточных алгоритмических
задач это прямой рецепт: один вызов, несколько кандидатов, тесты. Но
CodeMonkeys \cite{ehrlich2025codemonkeysscaling} показывает цену
многошагового сэмплирования на SWE-bench: 57{,}4\,\% при бюджете около
\$2\,300. Выбор между агентом и сэмплированием зависит от задачи, и его
надо делать явно.

\subsection{Переиспользование планов и опыта}

Шестая группа амортизирует расход между задачами. Agentic Plan Caching
\cite{zhang2025agenticplan} извлекает шаблоны планов из завершённых
прогонов и адаптирует их к похожим задачам: стоимость $-50{,}31\,\%$,
latency $-27{,}28\,\%$ при сохранении качества.

Для 100 скрытых задач эффект ограничен: задачи могут быть разными, а
холодный старт не даёт истории. Переиспользование оправдано, если в
прогоне есть семейства похожих задач, повторные попытки или общий набор
инструментов. В остальных случаях это дополнительная сложность без
гарантированного выигрыша.

\subsection{Неудачные задачи и распределение бюджета}

Седьмая группа улучшает знаменатель \eqref{eq:metric}: не даёт тратить
токены на задачи, которые не будут решены.

Ранний отказ в каскаде \cite{zellinger2025costsaving} даёт $-13\,\%$
стоимости и $-5\,\%$ ошибок. HERA \cite{luo2026heraharness} повышает
abstention accuracy до 83{,}3\,\% и переносит навык на другие модели.
BAVT \cite{li2026spendreason} при жёстком бюджете обходит baseline,
которому дали в четыре раза больше ресурса. Consequence-Aware Allocation
\cite{he2026allerrors} приоритизирует задачи по последствиям ошибки и
лучше решает высокоценные задачи при фиксированном compute. ROI-Reasoning
\cite{zhao2026roireasoning} формулирует бюджет как задачу о рюкзаке и
даёт до $+417\,\%$ относительного score при жёстком бюджете. ContextCurator
\cite{li2026escapingcontext} отдельной лёгкой политикой снижает токены на
8{,}8\,\% на WebArena и в 8 раз на DeepSearch.

Но отказ — не бесплатный: ложный отказ теряет баллы. В нашей постановке
невыполнимые задачи составляют 15--25\,\%, и их распознавание уже
исследовано в обзоре 0004. Здесь важен общий принцип: останавливать задачу
по калиброванному сигналу неуспеха, а не по фиксированному числу шагов.
Сигнал успеха можно получить по согласованности нескольких прогонов, но
сами прогоны стоят токенов, поэтому в бюджетном режиме нужен более дешёвый
признак.

\subsection{Hybrid-модели и обзоры}

HRBench \cite{ning2026hrbenchbenchmarkin} — наиболее близкая к нашей
постановке работа по выбору режима: шесть hybrid-моделей, пять бенчмарков,
12 настроек переключения. Prompt-based выбор режима даёт хороший
token--accuracy trade-off, внешний routing — более стабильное снижение
стоимости, speculative-исполнение обычно повышает и точность, и расход.
Обзоры \cite{alomrani2025reasoningbudget,sui2025stopoverthinking}
систематизируют управление test-time compute: фиксированный compute
неоптимален, нужна адаптация по трудности или уверенности.

\section{Сравнение}

Таблица ниже сводит включённые работы. Столбец «Результат» содержит число
из аннотации; оно не всегда переведено в нашу метрику.

{\footnotesize
\setlength{\tabcolsep}{3pt}
\begin{longtable}{@{}>{\raggedright\arraybackslash}p{2.6cm}>{\raggedright\arraybackslash}p{3.2cm}>{\raggedright\arraybackslash}p{3.2cm}>{\raggedright\arraybackslash}p{3.4cm}>{\raggedright\arraybackslash}p{3.0cm}@{}}
\toprule
Работа & Подход & Данные и метрика & Результат & Ограничение \\
\midrule
\endfirsthead
\toprule
Работа & Подход & Данные и метрика & Результат & Ограничение \\
\midrule
\endhead
\bottomrule
\endfoot
S01 Cost-of-Pass \cite{erol2025costpass} & стоимость правильного ответа & 3 класса задач, деньги & лёгкие/крупные/reasoning модели выгодны в разных классах; voting редко оправдан & доллары, не веса токенов \\
S02 Efficient Agents \cite{wang2025efficientagents} & упрощение архитектуры & GAIA, cost-of-pass & 96{,}7\,\% качества при $-\$0{,}17$; cost-of-pass $-28{,}4\,\%$ & один бенчмарк \\
S03 How Agents Spend \cite{bai2026aiagents} & анализ траекторий & SWE-bench, 8 моделей & 1000× токенов; разброс 30×; плато точности; трудность слабо предсказывает расход & сырые токены \\
S04 Tokenomics \cite{salim2026tokenomicsquantify} & разложение по фазам & ChatDev, 30 задач & code review 59{,}4\,\%; вход 53{,}9\,\% & один фреймворк \\
S05 Token Reduction \cite{weinberger2026tokenreduction} & сжатие вывода & кодовые задачи, Claude Code & $-38{,}4\,\%$ токенов, но $+6{,}8\,\%$ стоимости; $r=0{,}15$ & один провайдер \\
S06 OckBench \cite{du2025ockbenchmeasuring} & бенчмарк & рассуждение и код & разброс длины до 5× при близкой точности & без приёмов \\
S07 Harness Effect \cite{ali2026harnesseffect} & смена оркестрации & 22 задачи, 6 моделей & $-38\,\%$ токенов, $-41\,\%$ стоимости & малая выборка \\
S08 Local-Splitter \cite{agyemang2026localsplitter} & 7 тактик & 4 класса нагрузок & 45--79\,\% экономии; RAG 51\,\% & локальный триаж \\
S09 TokenEconomizer \cite{li2026tokeneconomizerexp} & внешний контроллер & последовательности задач & $-18{,}4$--$30{,}9\,\%$ стоимости & чёрный ящик \\
S10 TALE \cite{han2024tokenbudget} & бюджет в промпте & математика, QA & короче CoT при небольшой потере & подбор бюджета \\
S11 Token Complexity \cite{lee2025wellllms} & порог токенов & рассуждение & универсальный trade-off; нужна адаптация & нет агентов \\
S12 CARE \cite{he2026thinkor} & reward за длину & математика, код & $+4\,\%$ Pass@1 при $-37\,\%$ длины & обучение \\
S13 Dynamic Early Exit \cite{yang2025dynamicearly} & ранний выход & 10 бенчмарков, 11 моделей & $-19{,}1$--$80{,}1\,\%$ длины, $+0{,}3$--$5{,}0\,\%$ точности & нужны пробы \\
S14 ESTAR \cite{wang2026estarearly} & стоп-сигнал & 4 набора & 3{,}7× короче при 74{,}9 vs 74{,}2\,\% & обучение \\
S15 DTSR \cite{xiang2026thinkingenough} & проверка достаточности & Qwen3 & $-28{,}9$--$34{,}9\,\%$ длины & самооценка \\
S16 Budget Guidance \cite{li2025steeringllm} & предиктор длины & MATH-500 & 63\,\% токенов при конкурентной точности & математика \\
S17 L1 \cite{aggarwal2025l1controlling} & RL с длиной & рассуждение & плавный trade-off & обучение \\
S18 s1 \cite{muennighoff2025s1simple} & budget forcing & MATH, AIME & $+27\,\%$ над o1-preview; 50→57\,\% & математика \\
S19 DART \cite{lee2026dartdraft} & два черновика & математика, код & $-32$--$73\,\%$ thinking-токенов; $+9$/$+22{,}5$ п.п. & два черновика \\
S20 Ares \cite{yang2026aresadaptive} & роутер на шаг & TAU, BrowseComp, WebArena & $-52{,}7\,\%$ reasoning-токенов & обучение \\
S21 Confidence Training \cite{hosseini2026learningstop} & confidence как цель & 4 семейства моделей & $-25\,\%$ токенов при той же точности & обучение \\
S22 Danger of Overthinking \cite{cuadron2025dangeroverthinking} & анализ агентов & SWE-bench, 4018 траекторий & $+30\,\%$ успеха и $-43\,\%$ стоимости при выборе & нужны прогоны \\
S23 Harmful Overthinking \cite{caldarella2026thinkingpast} & остановка на префиксе & рассуждение & до $+21\,\%$ точности; early stop не лечит вредное & не агенты \\
S24 When More Thinking Hurts \cite{zhou2026thinkinghurts} & маргинальная польза & рассуждение & отдача падает; оптимум зависит от трудности & не агенты \\
S25 Reasoning Steps in Code \cite{xue2025empiricalstudy} & число шагов & BigCodeBench, 6 моделей & умеренное сокращение сохраняет стандарт, не трудное & 100 задач \\
S26 HERO \cite{li2026doingtokens} & иерархический RL & SWE-bench & лучший trade-off resolve/токены & обучение \\
S27 FrugalGPT \cite{chen2023frugalgptuse} & каскад & массовые запросы & качество GPT-4 при $-98\,\%$ стоимости & не агенты \\
S28 Early Abstention \cite{zellinger2025costsaving} & каскад с отказом & 6 бенчмарков & $-13\,\%$ стоимости, $-5\,\%$ ошибок, $+4{,}1\,\%$ отказов & QA \\
S29 Mixture of Thoughts \cite{yue2023largelanguage} & каскад по согласию & 6 наборов & качество GPT-4 за 40\,\% цены & старые модели \\
S30 STEER \cite{lee2025confidenceguided} & пошаговый routing & математика, QA, план & $+20\,\%$ точности при $-48\,\%$ FLOPs & FLOPs \\
S31 Cascading for Code \cite{chen2024modelcascading} & каскад + self-test & генерация кода & $-26\,\%$ стоимости, до $-70\,\%$ & code completion \\
S32 Agentic Routing \cite{zhang2026budgetaware} & выбор модели на шаг & агентские задачи & лучше cost--success frontier & обучение \\
S33 Compute-Optimal \cite{snell2024scalingllm} & аллокация compute & математика & 4× эффективнее best-of-N; 14× модель & FLOPs \\
S34 CodeMonkeys \cite{ehrlich2025codemonkeysscaling} & много траекторий & SWE-bench Verified & 57{,}4\,\% при $\approx\$2300$ & дорого \\
S35 Independent Sampling \cite{dong2026independentsamplin} & $k$-shot vs агент & 216 задач Codeforces & $k$-shot лучше по точности/стоимости & алгоритмические задачи \\
S36 Adaptive-Consistency \cite{aggarwal2023lets} & остановка голосования & 17 наборов & бюджет $-7{,}9×$ при потере $<0{,}1\,\%$ & только сэмплы \\
S37 LLMLingua \cite{jiang2023llmlinguacompressi} & сжатие промпта & 4 набора & до 20× сжатия & вход \\
S38 FOCUS \cite{dixit2026focustraining} & причинное сжатие & агентские бенчмарки & $-48\,\%$ контекста, $+8{,}9$ п.п. & не код \\
S39 Latent Tool \cite{liu2026toolsneed} & латентное сжатие & web QA & $-70$--$89\,\%$ токенов & обучение \\
S40 DTOC \cite{chaturvedi2026dtocdynamic} & обратимое сжатие & DeepSWE & cost per solved $-3$--$3{,}5×$ & зависит от модели \\
S41 ACM \cite{li2026acmagentic} & инструменты контекста & поиск и код & ниже пиковый контекст & демонстрации \\
S42 CAT \cite{liu2025contexttool} & контекст-инструмент & SWE-bench Verified & 57{,}6\,\% solved & обучение \\
S43 Context Folding \cite{sun2025scalinglong} & свёртка & research, SWE & 10× меньше активный контекст & RL \\
S44 ICLR \cite{wang2026agentsforget} & забывание reasoning & 260 задач & $-25{,}5\,\%$ входа, $-14{,}4\,\%$ выхода & не код \\
S45 SWE-MeM \cite{gao2026swemem} & память + GRPO & SWE-bench Verified & 43{,}4/60{,}2\,\% для 4B/30B & обучение \\
S46 TACO \cite{ren2026selfevolving} & сжатие наблюдений & TerminalBench, SWE-bench Lite & $+1$--$4\,\%$ точности, ниже расход & терминал \\
S47 TraceLite \cite{xiao2026traceliteefficient} & изображения + дубли & SWE-bench, DeepSWE & 78{,}25\,\% resolve, $-33{,}05\,\%$ входа & рендер \\
S48 Measure Before Manage \cite{chen2026measurebefore} & анализ памяти & 55 траекторий & равный бюджет $\neq$ равный контекст & мало данных \\
S49 Harness Design \cite{fan2026empiricalstudy} & компоненты harness & 176 настроек & elision до суммаризации лучше; planning экономит & конкретный harness \\
S50 AGENTS.md \cite{lulla2026impactagents} & инструкции репозитория & 10 репозиториев, 124 PR & $-16{,}58\,\%$ выходных токенов, $-28{,}64\,\%$ runtime & наблюдательное \\
S51 FastContext \cite{zhang2026fastcontexttrainin} & отдельный исследователь & SWE-bench & $+5{,}5$ п.п. resolve, $-60\,\%$ токенов & обучение \\
S52 Agentless \cite{xia2024agentlessdemystify} & 3 фазы & SWE-bench Lite & 32{,}00\,\% при \$0{,}70 & модели 2024 \\
S53 Tool Architecture \cite{xu2026devilinterface} & 6 архитектур & 11\,700 траекторий & $-41{,}6\,\%$ шагов, $-56{,}3\,\%$ токенов & конкретные акторы \\
S54 Steer, Don't Solve \cite{gandhi2026steerdon} & дешёвый критик & SWE-bench, 6 агентов & $+16{,}0$/$+14{,}4$ п.п.; \$0{,}07→\$0{,}03 & обучение \\
S55 ContextPipe \cite{xu2026contextpipedatabas} & сборка контекста & SWE-bench Pro & $-31\,\%$ токенов, $-23\,\%$ вызовов & один срез \\
S56 Plan Caching \cite{zhang2025agenticplan} & кеш планов & агентские приложения & $-50{,}31\,\%$ стоимости & похожие задачи \\
S57 BAVT \cite{li2026spendreason} & budget-aware search & 4 набора QA & лучше baseline с 4× ресурсом & QA \\
S58 Consequence-Aware \cite{he2026allerrors} & приоритет по последствиям & SWE-bench Lite & лучший успех на высокоценных задачах & нужен предиктор \\
S59 HERA \cite{luo2026heraharness} & коэволюция harness & выполнимые и невыполнимые & abstention 61{,}7→83{,}3\,\%; $-85\,\%$ стоимости & синтетика \\
S60 ROI-Reasoning \cite{zhao2026roireasoning} & бюджет как рюкзак & математика, MCQ, код & до $+417\,\%$ score при жёстком бюджете & обучение \\
S61 ContextCurator \cite{li2026escapingcontext} & отдельная политика & WebArena, DeepSearch & $-8{,}8\,\%$ токенов; 8× на DeepSearch & web/search \\
S62 HRBench \cite{ning2026hrbenchbenchmarkin} & переключение режима & 6 hybrid-моделей & prompt-based; routing стабильнее & бенчмарк \\
S63 Reasoning on a Budget \cite{alomrani2025reasoningbudget} & обзор & reasoning-модели & L1-контроль и L2-адаптивность & обзор \\
S64 Stop Overthinking \cite{sui2025stopoverthinking} & обзор & reasoning-модели & таксономия overthinking & обзор \\
\end{longtable}
}

\section{Пробел и вывод}

\subsection{Что установлено}

Из 64 работ складывается устойчивая картина.

\begin{enumerate}
  \item \textbf{Выход дороже входа.} В нашей метрике один выходной токен
        стоит десять входных. Поэтому приёмы, меняющие рассуждение и число
        выходных шагов, дают больший рычаг, чем сжатие промпта. Это
        подтверждается и числами: сжатие вывода инструментов на 38{,}4\,\%
        может увеличить стоимость \cite{weinberger2026tokenreduction}, а
        смена режима рассуждения снижает thinking-токены на 32--73\,\%
        \cite{lee2026dartdraft}.
  \item \textbf{Сокращение токенов не равно сокращению стоимости.} Сжатие
        меняет траекторию: агент дольше ищет, чаще проверяет, делает лишние
        шаги. Мерить надо cost per solved task, а не токены
        \cite{weinberger2026tokenreduction,bai2026aiagents}.
  \item \textbf{Адаптивность важнее среднего.} Единый бюджет рассуждения
        проигрывает: у задач разная токенная сложность
        \cite{lee2025wellllms}, а оптимальная длина зависит от трудности
        \cite{zhou2026thinkinghurts}. Выигрывают методы, которые оценивают
        сложность или уверенность и меняют бюджет на задаче или шаге
        \cite{han2024tokenbudget,yang2025dynamicearly,yang2026aresadaptive,lee2026dartdraft}.
  \item \textbf{Число шагов — отдельный и большой рычаг.} CodeAct-подобный
        интерфейс, отдельный исследователь репозитория и дешёвый критик
        дают 40--60\,\% экономии \cite{xu2026devilinterface,zhang2026fastcontexttrainin,gandhi2026steerdon}.
        Для самодостаточных задач независимое сэмплирование может быть
        дешевле агента \cite{dong2026independentsamplin}.
  \item \textbf{Неудачная задача — чистая стоимость.} Ранний отказ и
        abstention улучшают знаменатель \cite{zellinger2025costsaving,luo2026heraharness},
        но ложный отказ теряет баллы; нужен калиброванный сигнал.
  \item \textbf{Контекст полезен, пока не начинает менять траекторию.}
        Свежие наблюдения лучше хранить точно, старые — сжимать или
        выгружать; обратимость помогает не всегда
        \cite{fan2026empiricalstudy,chen2026measurebefore}.
\end{enumerate}

\subsection{Пробел}

Ни одна включённая работа не измеряет нашу метрику напрямую: четыре модели
Qwen, 100 кодовых задач, общий бюджет, один час, пять параллельных сессий,
$C=O+\lceil I/10\rceil$. Почти все работы используют проприетарные модели,
доллары, FLOPs или сырые токены; кодовые агенты чаще всего оцениваются на
SWE-bench, а не на коротких задачах с баллами. Поэтому числа обзора — это
ориентир, а не готовый рецепт.

Открытые вопросы, на которые обзор не отвечает:

\begin{itemize}
  \item где проходит порог переключения на \code{-noreason} для Qwen3.8/3.6
        на наших задачах;
  \item сколько входных токенов выгодно потратить, чтобы сэкономить один
        выходной: при весе 1/10 это зависит от задачи и от того, сколько
        шагов экономит контекст;
  \item насколько ранний отказ улучшает \eqref{eq:metric} при 15--25\,\%
        невыполнимых задач и как не потерять выполнимые;
  \item сколько стоит разброс: средняя стоимость может быть низкой, а
        хвост — съедать бюджет.
\end{itemize}

\subsection{Что делать}

Портфель приёмов, упорядоченный по силе доказательств и переносимости:

\begin{enumerate}
  \item \textbf{Маршрутизация режима рассуждения.} Начать с
        \code{-noreason}; повышать режим при расхождении двух дешёвых
        черновиков, провале тестов или признаках неуверенности. Ожидаемое
        снижение — 30--70\,\% выходных токенов; доказательства:
        \cite{lee2026dartdraft,yang2026aresadaptive,ning2026hrbenchbenchmarkin,chen2023frugalgptuse}.
  \item \textbf{Бюджет длины с ранней остановкой.} Короткий бюджет по
        умолчанию, повышение на трудных задачах, остановка по достаточности
        или первому стабильному ответу. Ожидаемое снижение — 20--50\,\%
        длины; доказательства:
        \cite{han2024tokenbudget,yang2025dynamicearly,wang2026estarearly,li2025steeringllm,hosseini2026learningstop}.
  \item \textbf{Сокращение шагов.} Интерфейс действий в стиле CodeAct,
        план перед редактированием, дешёвый критик, отдельный исследователь
        репозитория. Ожидаемое снижение — 40--60\,\% шагов и токенов;
        доказательства:
        \cite{xu2026devilinterface,zhang2026fastcontexttrainin,gandhi2026steerdon,xia2024agentlessdemystify}.
  \item \textbf{Политика контекста.} Свежие наблюдения — точно, старые —
        сжимать или выгружать, восстановление по требованию; входные
        инструкции и релевантный код давать сразу, потому что вход дешёвый.
        Снижение — 20--50\,\% входа и меньше переполнений; доказательства:
        \cite{liu2025contexttool,sun2025scalinglong,xiao2026traceliteefficient,fan2026empiricalstudy,lulla2026impactagents}.
  \item \textbf{Ранний отказ.} Калиброванный сигнал неуспеха после
        нескольких непродуктивных шагов; отдельно — детектор невыполнимых
        задач из обзора 0004. Снижение — 10--20\,\% стоимости за счёт
        неудачных задач; доказательства:
        \cite{zellinger2025costsaving,luo2026heraharness,he2026allerrors}.
  \item \textbf{Переиспользование планов.} Только для повторяющихся
        семейств задач и повторных попыток; доказательство —
        \cite{zhang2025agenticplan}.
\end{enumerate}

Чего не делать в первую очередь:

\begin{itemize}
  \item кеш провайдера и speculative decoding — не меняют число токенов в
        нашей метрике;
  \item сжатие промпта само по себе — вход в десять раз дешевле выхода,
        а эффект на успех не гарантирован \cite{weinberger2026tokenreduction};
  \item увеличение числа сэмплов без селектора — числитель растёт быстрее
        знаменателя \cite{ehrlich2025codemonkeysscaling};
  \item многоагентные схемы — добавляют вызовы и передачу контекста, пока
        не доказано обратное \cite{wang2025efficientagents}.
\end{itemize}

\subsection{Эксперимент}

До спецификации изменения нужно поставить эксперимент в
\code{research/experiments/}. Минимальный дизайн:

\begin{itemize}
  \item фиксированный набор кодовых задач, стратифицированный по трудности
        и включающий невыполнимые задачи;
  \item базовая линия — текущий ReAct с полным рассуждением и полным
        контекстом;
  \item варианты — по одному приёму за раз, затем лучший портфель;
  \item метрики — $C$ на задачу, $C$ на решённую задачу, доля решённых,
        медиана и 90-й процентиль стоимости, число шагов;
  \item три повтора на конфигурацию; парное сравнение по задачам;
  \item отдельная проверка, что экономия не достигнута ценой падения
        успеха.
\end{itemize}

Если приём даёт устойчивый прирост на фиксированном бюджете, по нему
пишется спецификация изменения в \code{specification/changes/}. До этого
обзор остаётся исследованием, а не контрактом поведения.

\section{Ограничения обзора}

\begin{itemize}
  \item Две базы; Semantic Scholar и OpenAlex недоступны, Crossref плохо
        ранжирует CS-препринты, поэтому почти все включённые работы найдены
        в arXiv.
  \item Просматривались верхние 20 записей на запрос, а не вся выдача.
  \item Данные извлечены из аннотаций; полные тексты проверялись
        выборочно. Числа могут отличаться от итоговых таблиц статей.
  \item Метрики несопоставимы: длина CoT, входные и выходные токены,
        доллары, FLOPs, число шагов, cost per solved task. Мета-анализ
        невозможен.
  \item Почти нет работ, которые меряют именно $O+I/10$ на кодовом агенте
        с четырьмя моделями Qwen и общим бюджетом; перенос чисел — оценка,
        а не факт.
  \item Большинство работ 2025--2026 — препринты, часть не прошла
        рецензирование.
  \item Обзор не оценивает публикационное смещение и не воспроизводит
        замеры.
\end{itemize}

\bibliographystyle{plain}
\bibliography{refs}

\end{document}
